% MRR, panel (a): what the reciprocal rank of a single query is, and how the
% mean over queries is formed. Each column is one query spectrum; the gallery
% is ranked by cosine similarity and the amber cell is the true molecule.
% Requires: shared/notation.tex (fig styles, amber), amssymb (\checkmark).
\begin{tikzpicture}[x=1cm,y=1cm,fig,>={Latex[length=2.6mm]}]

  \def\w{1.7}    % column width
  \def\h{0.62}   % cell height
  \def\gap{2.25} % column pitch
  \def\n{6}      % candidates shown

  % rank axis (shared by all three columns)
  \foreach \i in {1,...,\n}
    \node[font=\scriptsize,black!55,anchor=east] at (-0.18,{-(\i-0.5)*\h}) {\i};
  \node[font=\scriptsize,black!55,rotate=90,anchor=south]
       at (-0.62,{-0.5*\n*\h}) {rank};
  \draw[black!40,fig thin,->] (-0.95,-0.12) -- (-0.95,{-\n*\h+0.12});

  % three example queries, with the true molecule at a different rank each
  \foreach \q/\hit/\rr [count=\k from 0] in {1/1/1, 2/3/{\tfrac{1}{3}}, 3/2/{\tfrac{1}{2}}}{
    \pgfmathsetmacro{\x}{\k*\gap}
    \node[font=\scriptsize,amber!30!black,anchor=south] at ({\x+0.5*\w},0.52) {query \q};
    \draw[amber!70!black,fig line,->] ({\x+0.5*\w},0.44) -- ({\x+0.5*\w},0.08);
    \foreach \i in {1,...,\n}{
      \ifnum\i=\hit
        \fill[amber!45,draw=amber!70!black,line width=0.5pt]
             ({\x},{-(\i-1)*\h}) rectangle ++({\w},{-\h});
        \node[font=\scriptsize,amber!25!black] at ({\x+0.5*\w},{-(\i-0.5)*\h}) {$\checkmark$ true};
      \else
        \fill[RoyalBlue!10,draw=black!35,line width=0.4pt]
             ({\x},{-(\i-1)*\h}) rectangle ++({\w},{-\h});
      \fi
    }
    \node[font=\small,anchor=north] at ({\x+0.5*\w},{-\n*\h-0.20})
         {$\mathrm{RR}=\rr$};
  }

  % candidate-library caption
  \node[font=\scriptsize,black!55,anchor=south west] at (-0.95,1.12)
       {gallery ranked by cosine similarity};

  % the mean over queries
  \node[font=\small,anchor=north] at ({\gap+0.5*\w},{-\n*\h-1.05})
       {$\displaystyle\mathrm{MRR}=\frac{1}{N}\sum_{i=1}^{N}\frac{1}{r_i}
         =\frac{1+\tfrac{1}{3}+\tfrac{1}{2}}{3}=0.61$};

\end{tikzpicture}
